\section{Structure of the Course}

The course is organized in three main parts, followed by appendices.

Part I develops the foundations. It explains how field theory grows out of
mechanics, introduces the mathematical language of fields, and reviews the
spacetime formulation needed for relativistic theories.

Part II develops the Lagrangian and Hamiltonian formalisms for fields. The
central objects are the action functional, the Lagrangian density, the
Euler--Lagrange equations, symmetries, Noether currents, conserved charges, and
canonical phase space.

Part III studies fundamental classical fields: real and complex scalar fields,
the electromagnetic field, non-Abelian gauge fields, and spinor fields. These
examples are chosen because they form the classical backbone of many later
quantum and gauge-theoretic constructions.

The appendices collect conventions and technical tools: variational calculus,
Lie groups and Lie algebras, differential forms, Green's functions, and problem
sets. They are meant to support the main text without interrupting the main
line of development.

\begin{remarkbox}
The book should be read actively. A field equation should be differentiated,
varied, dimensionally checked, solved in simple cases, and interpreted. The
notes are designed to grow by this kind of work.
\end{remarkbox}
